\documentclass[12 pt, letterpaper]{hmcpset}
\usepackage{hw-preamble}

\name{Jayden Thadani}
\class{MATH 135 --- Complex Analysis}
\assignment{Homework 2}
\duedate{12th September 2024}

\begin{document}

Collaborators: Kai Mawhinney, Justin Son

\begin{problem}[16.]
	Determine the radius of convergence of the series $\sum_{n = 1}^{\infty} a_{n} z^{n}$ when:
	\begin{enumerate}
		\item $a_{n} = (\log n)^{2}$ 
		\item $a_{n} = (n!)^{3}/(3n)!$.
	\end{enumerate}
\end{problem}

\pagebreak

\begin{problem}[17.]
	Show that if $\{ a_{n} \}_{n = 0}^{\infty}$ is a sequence of non-zero complex numbers such that
	\[
		\lim_{n \to \infty} \frac{\lvert a_{n + 1} \rvert}{\lvert a_{n} \rvert} = L,
	\]
	then
	\[
		\lim_{n \to \infty} \lvert a_{n} \rvert^{1/n} = L.
	\]
\end{problem}

\begin{solution}
	Suppose we have that $\lim_{n \to \infty} \lvert a_{n + 1} \rvert/\lvert a_{n} \rvert = L$. Then for every, $\varepsilon$, we must have some $N$ such that for every $n > N$
	 \[
		\left \lvert \frac{\lvert a_{n + 1} \rvert}{\lvert a_{n} \rvert} - L \right \rvert < \varepsilon
	\]
	Thus,
	\[
		(L - \varepsilon) \lvert a_{n} \rvert < \lvert a_{n + 1} \rvert < (L + \varepsilon) \lvert a_{n} \rvert. 
	\]
	Then by induction, we can see that for any $k$, we must have
	\[
		(L - \varepsilon)^{k} \lvert a_{n} \rvert < \lvert a_{n + k} \rvert < (L + \varepsilon)^{k} \lvert a_{n} \rvert.
	\]
	Taking $(n + k)$th roots, we get
	\[
		(L - \varepsilon)^{\frac{k}{n + k}} \lvert a_{n} \rvert^{\frac{1}{n + k}} < \lvert a_{n + k} \rvert^{\frac{1}{n + k}} < (L + \varepsilon)^{\frac{k}{n + k}} \lvert a_{n} \rvert^{\frac{1}{n + k}}
	\]
	Choose 
\end{solution}

\pagebreak

\begin{problem}[18.]
	Let $f$ be a power series centred at the origin. Prove that $f$ has a power series expansion around any point in its disc of convergence.
\end{problem}

\pagebreak

\begin{problem}[20.]
	Expand $(1 - z)^{-m}$ in powers of $z$. Here $m$ is a fixed positive integer. Also, show that if
	\[
		(1 - z)^{-m} = \sum_{n = 0}^{\infty} a_{n} z^{n},
	\]
	then one obtains the following asymptotic relation for the coefficients:
	\[
		a_{n} \sim \frac{1}{(m - 1)!} n^{m - 1} \text{ as } n \to \infty.
	\]
\end{problem}

\pagebreak

\begin{problem}[25.]
	The next three calculations provide some insight into Cauchy's theorem, which we treat in the next chapter.
	\begin{enumerate}
		\item Evaluate the integrals
			\[
				\int_{\gamma} z^{n} \, dz 
			\]
			for all integers $n$. Here $\gamma$ is any circle centred at the origin with the positive orientation.
		\item Same question as before, but with $\gamma$ any circle not containing the origin.
		\item Show that if $\lvert a \rvert < r < \lvert b \rvert$, then
			\[
				\int_{\gamma} \frac{1}{(z - a)(z - b)} \, dz = \frac{2 \pi i}{a - b}
			\]
			where $\gamma$ denotes the circle centred at the origin, of radius $r$, with the positive orientation.
	\end{enumerate}
\end{problem}

\begin{solution}
	\begin{enumerate}
		\item 
			\[
				\boxed{
					\int_{\gamma} z^{n} \, dz = \begin{cases}
						0 & n \neq -1 \\
						2 \pi i r^{n} & n = -1.
					\end{cases}
				}
			\]

			Consider the parametrisation $z : t \to r e^{i t}$ for $t \in [0, 2 \pi]$. Then
			\[
				\begin{split}
					\int_{\gamma} z^{n} \, dz & = \int_{0}^{2 \pi} (r e^{i t})^{n} ir e^{it} \, dt \\
								  & = i r^{n} \int_{0}^{2 \pi} e^{(n + 1) i t} \, dt.
				\end{split}
			\]
			If $n = -1$, then $e^{(n + 1) i t} = 1$ and the integral is just $2 \pi$, giving us
			\[
				\int_{\gamma} z^{n} \, dz = 2 \pi i r^{n}.
			\]
			If $n \neq -1$, we can use the fundamental theorem of calculus since $\frac{1}{i(n + 1)} e^{(n + 1) i t}$ is a primitive of $e^{(n + 1) i t}$.
	\end{enumerate}
\end{solution}

\pagebreak

\begin{problem}[5.]
	Suppose $f$ is continuously \emph {complex} differentiable on $\Omega$, and $T \subset \Omega$ is a triangle whose interior is also contained in $\Omega$. Apply Green's theorem to show that
	\[
		\int_{T} f(z) \, dz = 0.
	\]
\end{problem}

\pagebreak

\begin{problem}[6.]
	Let $\Omega$ be an open subset of $\mathbb{C}$ and let $T \subset \Omega$ be a triangle whose interior is also contained in $\Omega$. Suppose that $f$ is a function holomorphic in $\Omega$ except possibly at a point $w$ inside $T$. Prove that if $f$ is bounded near $w$, then
	\[
		\int_{T} f(z) \, dz = 0.
	\]
\end{problem}

\end{document}
